\documentclass{homework}
\author{Tashfeen, Ahmad}
\class{CSCI 6213: Data Science Fundamentals}
\title{Homework 6}
\address{Oklahoma City University, Petree College of Arts \& Sciences, Computer Science}

\begin{document} \maketitle
You're \textit{not allowed} to use any libraries other than Numpy or Matplotlib.

Let $x_{1}, \ldots, x_n$ be a set of instances in $\R^p$ just like before and $c_1, \ldots, c_K$ be a
set of $K$ cluster centroids. Note that $c_i$ are also in $\R^p$. Let $\mathbb{I}_{ij} \in \{0, 1\}$ be a variable such that $\mathbb{I}_{ij} = 1$ if $x_i$ is assigned to cluster centered at $c_j$  and $\mathbb{I}_{ij} = 0$ otherwise. In layman's terms, $\mathbb{I}_{ij}$ indicates whether $x_i$ is in the $j^\text{th}$ cluster. The $K$--Means Clustering algorithm minimises the following objective function,
\[
  \sum_{i=1}^n \sum_{j=1}^K \mathbb{I}_{ij} [(x_i - c_j)^T(x_i - c_j)]
\]
This objective function is essentially a measure of how far are the members of each cluster from their respective centroids or means. The $K$--Means Clustering algorithm is given bellow,

\begin{enumerate}
  \item Initialise centroids or means $c_j$ for $1 \leq j \leq K$ randomly in $\R^p$ or,
        \[
          c_j \leftarrow x_j
        \]
  \item Assign all $x_i$ to the closest centroid $c_u$. I. e., for some $u$ and all $j$,
        \[
          \mathbf{if}\quad [(x_i - c_u)^T(x_i - c_u)]<[(x_i - c_j)^T(x_i - c_j)]\quad \mathbf{then} \quad \mathbb{I}_{iu}\leftarrow 1
        \]
  \item Adjust the centroids or means according the new assignment in the last step,
        \[
          c_j \leftarrow \sum_{i=1}^{n}\mathbb{I}_{ij}\paren{\frac{x_i}{n}}
        \]
        Note that this is just assigning $c_j$ to the new mean of the $j^\text{th}$ cluster.
  \item Repeat from (2) till $c_j$ stop changing or the change is less than some $\ep$.
\end{enumerate}
You may also find this algorithm (Algorithm 12.2) on page 523 of your textbook (Hatie et al.).
\lstinputlisting[
  language={Python},
  caption={Python code to generate synthetic two dimensional data clusters. \href{https://tashfeen.org/raw/share/dsf/kmeans.py}{Download} this file.},
  label=code]
{./code/kmeans.py}

\question Listing \ref{code} plots the clusters shown in figure \ref{clusters}. Give your best guesses for the optimal centroids $c_j\in \R^2$ and $K=4$.

\question Implement the $K$--Means Clustering algorithm for the data in listing \ref{code}. Run the algorithm for $K \in \{1, 2, 3, 4\}$ and $\ep = 10^{-5}$. Give your final converged centroids plotted over the clustered data. Plot each cluster with a distinct colour\footnote{use different shades of grey if you have difficulty seeing colour}. At the end you should have four converged (with different coloured clusters) plots for each value of $K$.


\question Download the file \href{https://myokcuedu-my.sharepoint.com/:u:/g/personal/tashfeen_okcu_edu/EeqU1CH36W5Bp0YRgiRXxxUBuhzdI_JPec6CBlxsMBCtgA}{\texttt{t10k-images-idx3-byte}}. This is the binary file containing 1000 images of the test set from the \href{http://yann.lecun.com/exdb/mnist/}{the MNIST database of handwritten digits}. You can refer to their website for more information about the data and its organisation. However, it is read and plotted for your inspection and further usage to implement the $K$--Means Clustering algorithm in listing \ref{mnsit}.

\lstinputlisting[
  language={Python},
  caption={Python code that reads in the binary MNSIT test set data and plots a few data points. Pressing the space-bar changes the plot upon running the file. \href{https://tashfeen.org/raw/share/dsf/mnsit.py}{Download} this file.},
  label=mnsit]
{./code/mnsit.py}

\begin{enumerate}
    \item Why do you think $K=10$ maybe a good choice for this data?
    \item What are the values of $n$ and $p$?
    \item What does each of the total $p$ features represent in this data?
    \item For $\ep = 3.5$, finish implementing the $K$--Means Clustering algorithm after the line indicated by the comment in listing \ref{mnsit}. The code in listing \ref{mnsit} should then plot your 10 approximately converged cluster centroids as 28 by 28 pixel images comprising of 2 rows and 5 columns. Give and discuss this plot. Why do you think the centroids $c_j$ look the way they do?
\end{enumerate}

\img<clusters>[0.7]{Unlabelled clusters of data.}{./media/clusters}

\end{document}
